% Part: sets-functions-relations
% Chapter: functions
% Section: kinds

\documentclass[../../../include/open-logic-section]{subfiles}

\begin{document}

\olfileid{sfr}{fun}{kin}
\olsection{Verschillende soorten functies}

\begin{explain}
Het helpt om de soorten functies die we vaak tegenkomen in groepen in
te delen.

We beginnen met functies waarbij ieder element van het codomein ook
echt als waarde voorkomt. Zo'n functie heet !!{surjective}. De
afbeelding in \olref{fig:surjective} laat dit zien.

\begin{figure}
  \olasset{assets/diagrams/surjective.tikz}
  \caption{Bij !!^a{surjective} functie komt ieder !!{element} van het
    codomein als waarde voor.}
  \ollabel{fig:surjective}
\end{figure}
\end{explain}

\begin{defn}[!!^{surjective} functie]
Een functie $f \colon A \rightarrow B$ heet \emph{!!{surjective}}
precies wanneer $B$ ook het bereik van~$f$ is. Dat betekent: voor elke
$y \in B$ is er minstens één $x \in A$ waarvoor~$f(x) = y$, of in
symbolen:
\[
  (\forall y \in B)(\exists x \in A)f(x) = y.
\]
Zo'n functie noemen we !!a{surjection} van $A$ naar $B$.
\end{defn}

\begin{explain}
Wil je laten zien dat $f$ !!a{surjection} is, dan moet je laten zien
dat ieder object in het codomein van $f$ de waarde $f(x)$ is bij een
invoer $x$.

Elke functie \emph{levert} !!a{surjection} op. Neem namelijk een
functie $f \colon A \to B$ en maak $f' \colon A \to \ran{f}$ met
$f'(x) = f(x)$. Het bereik $\ran{f}$ is \emph{per definitie}
$\Setabs{f(x) \in B}{x \in A}$. Daarom is deze functie $f'$ zeker
!!a{surjection}.
\end{explain}

\begin{explain}
Elke mogelijke invoer van een functie krijgt precies één uitvoer. Bij
sommige functies krijgen twee verschillende invoeren bovendien nooit
dezelfde uitvoer. Zo'n functie heet !!{injective}. De afbeelding in
\olref{fig:injective} laat dit zien.
\begin{figure}
  \olasset{assets/diagrams/injective.tikz}
  \caption{!!^a{injective} functie geeft nooit dezelfde waarde aan
    twee verschillende argumenten.}
  \ollabel{fig:injective}
\end{figure}
\end{explain}

\begin{defn}[!!^{injective} functie]
Een functie $f \colon A \rightarrow B$ heet \emph{!!{injective}}
precies wanneer er bij elke $y \in B$ hoogstens één $x \in A$ is
waarvoor~$f(x) = y$. Zo'n functie noemen we !!a{injection} van $A$
naar~$B$.
\end{defn}

\begin{explain}
Wil je laten zien dat $f$ !!a{injection} is, neem dan willekeurige
!!{element}s $x$ en $y$ uit het domein van $f$. Je moet bewijzen: als
$f(x)=f(y)$, dan $x=y$.
\end{explain}

\begin{ex}
De constante functie $f\colon \Nat \to \Nat$ met $f(x) = 1$ is niet
!!{injective} en ook niet !!{surjective}.

De identiteitsfunctie $f\colon \Nat \to \Nat$ met $f(x) = x$ is
!!{injective} én !!{surjective}.

De opvolgerfunctie $f \colon \Nat \to \Nat$ met $f(x) = x+1$ is
!!{injective}, maar niet !!{surjective}.
  
De functie $f \colon \Nat \to \Nat$ die als volgt is vastgelegd,
\[
  f(x) =
  \begin{cases}
    \frac{x}{2} & \text{als $x$ even is} \\
    \frac{x+1}{2} & \text{als $x$ oneven is.}
  \end{cases}
\]
is !!{surjective}, maar niet !!{injective}.
\end{ex}

\begin{explain}
Vaak hebben we een functie nodig die tegelijk !!{injective} en
!!{surjective} is. Zo'n functie heet !!{bijective}. Zij ziet eruit als
de functie in \olref{fig:bijective}. !!^{bijection}s heten ook wel
\emph{een-op-een-koppelingen}: ze koppelen ieder element van het
codomein aan precies één element van het domein.
\begin{figure}
  \olasset{assets/diagrams/bijective.tikz}
  \caption{!!^a{bijective} functie koppelt ieder element van het
    codomein aan precies één element van het domein.}
  \ollabel{fig:bijective}
\end{figure}
\end{explain}

\begin{defn}[!!^{bijection}]
Een functie $f \colon A \to B$ heet \emph{!!{bijective}} precies
wanneer zij zowel !!{surjective} als !!{injective} is. Zo'n functie
noemen we !!a{bijection} van $A$ naar~$B$ (of tussen $A$ en~$B$).
\end{defn}
\end{document}
